%%%%%%%%%%%%%%%%%%%%%%%%%%%%%%%%%%%%%%%%%%%%%%%%%%%%%%%%%%%%%%%%%%%%%%%%%
\begin{easyappendix}{Struktura repozytorium kodu i~instrukcja uruchomieniowa}
\label{app:kod}

Niniejszy załącznik zawiera techniczną dokumentację wytworzonego oprogramowania, opis struktury repozytorium oraz instrukcję reprodukcji wyników eksperymentów numerycznych i~optymalizacyjnych Bliźniaka Cyfrowego. Całość kodu źródłowego pracy została objęta systemem kontroli wersji Git i~jest dostępna w~repozytorium projektowym.

\section*{Wymagania środowiskowe i instalacja}

Projekt został zaimplementowany w~języku \textbf{Python~3.12} (z~pełną kompatybilnością z~wersjami 3.10--3.14). Zarządzanie zależnościami oparto o~środowisko wirtualne \texttt{venv} oraz zablokowane wersje bibliotek w~pliku \texttt{requirements.txt}.

\begin{lstlisting}[language=bash,
	caption={Kroki instalacji środowiska wirtualnego i~zależności},
	label={lst:instalacja}]
	# 1. Klonowanie repozytorium
	git clone https://github.com/rafal-kaczynski-raf-tech/praca-magisterska.git
	cd praca-magisterska

	# 2. Utworzenie i aktywacja wirtualnego srodowiska
	python3 -m venv .venv
	source .venv/bin/activate        # Linux / macOS
	# .venv\Scripts\activate         # Windows

	# 3. Instalacja scisle zdefiniowanych zaleznosci
	pip install --upgrade pip
	pip install -r requirements.txt
\end{lstlisting}

\section*{Główne zależności obliczeniowe i narzędziowe}
\begin{itemize}
	\item \textbf{NumPy} ($\ge 2.0.0$) -- wektoryzowane struktury danych, tablice próbek i~operacje całkowania trapezowego,
	\item \textbf{PySwarms} ($= 1.3.0$) -- zweryfikowana biblioteka akademicka kanonicznego algorytmu PSO (\texttt{GlobalBestPSO}),
	\item \textbf{SciPy} ($\ge 1.18.0$) -- funkcje statystyczne i~pomocnicze procedury naukowe,
	\item \textbf{Matplotlib} ($\ge 3.10.0$) -- generowanie publikacyjnych wykresów i~trajektorii czasowych,
	\item \textbf{ReportLab} ($\ge 4.0.0$) -- automatyczny generator formalnych raportów akademickich w~formacie PDF.
\end{itemize}

\section*{Struktura katalogów i modułów projektu}
Projekt podzielono na modularne pakiety zgodnie z~zasadami czystej architektury (ang.~\emph{Clean Architecture}):

\begin{itemize}
	\item \directory{src/} -- Rdzeń Bliźniaka Cyfrowego (silnik fizyczny i~sterowniki):
	\begin{itemize}
		\item \texttt{converter.py} -- model stanu przekształtnika podwyższającego (Boost) i~implementacja metody Eulera,
		\item \texttt{controller.py} -- implementacja regulatorów: histerezowego, klasycznego MF-BB oraz proponowanego Trend MF-BB z~ogranicznikiem prądu,
		\item \texttt{simulator.py} -- orkiestrator pętli czasowej sprzęgający fizykę obwodu z~dyskretnym taktowaniem sterownika,
		\item \texttt{filters.py} -- dyskretne filtry dolnoprzepustowe 1. rzędu Tustina ($T_s = 10~\mu\mathrm{s}$),
		\item \texttt{config.py} -- klasy konfiguracyjne parametrów obwodu i~symulacji oparte o~dekorator \texttt{@dataclass}.
	\end{itemize}
	\item \directory{optymalizacja/} -- Moduły optymalizacji rojem cząstek (PSO) i~analizy wrażliwości:
	\begin{itemize}
		\item \texttt{pyswarms\_optimizer.py} -- oficjalny adapter łączący symulator z~biblioteką \texttt{PySwarms},
		\item \texttt{cost\_functions.py} -- implementacja funkcji celu (MSE, IAE, ITAE, CurrentAware, transient ISE),
		\item \texttt{optymalizacja\_dwuetapowa\_2d.py} -- skrypt wykonawczy procedury dwutorowej 2D PSO,
		\item \texttt{badanie\_wrazliwosci.py} -- analiza odporności i~wrażliwości wag $\lambda_i$ oraz $w_{\mathrm{cross}}$,
		\item \texttt{test\_pyswarms\_porownanie.py} -- skrypt walidacji krzyżowej silnika PySwarms względem procedury referencyjnej.
	\end{itemize}
	\item \directory{tests/} -- Zautomatyzowany zestaw testów jednostkowych i~regresyjnych:
	\begin{itemize}
		\item \texttt{test\_regression.py} -- test regresyjny weryfikujący zgodność bit-w-bit z~zamrożonym baseline'em,
		\item \texttt{test\_filters.py} -- testy charakterystyk częstotliwościowych filtrów cyfrowych.
	\end{itemize}
\end{itemize}

\section*{Instrukcja uruchomienia kluczowych eksperymentów}

\begin{lstlisting}[language=bash,
	caption={Polecenia uruchomienia testow, symulacji oraz optymalizacji},
	label={lst:uruchomienie}]
	# 1. Uruchomienie testow jednostkowych i regresyjnych
	pytest tests/

	# 2. Uruchomienie pelnej symulacji porownawczej wzgledem modelu PSIM
	python porownanie_bb_estim_full_psim.py

	# 3. Wykonanie dwutorowej optymalizacji 2D z uzyciem PySwarms
	python -m optymalizacja.optymalizacja_dwuetapowa_2d

	# 4. Wykonanie pelnego badania wrazliwosci algorytmu PSO
	python -m optymalizacja.badanie_wrazliwosci
\end{lstlisting}

Wyniki symulacji numerycznych oraz wygenerowane raporty w~formacie PDF i~wykresy PNG zapisywane są bezpośrednio w~katalogu głównym projektu oraz podkatalogu \directory{optymalizacja/}.

\end{easyappendix}
	\caption{Wykorzystanie programu Kile do edycji tabeli}
	\label{rys:kile}
\end{figure}

Dedykowane edytory, takie jak \href{https://kile.sourceforge.io/}{Kile}, również często oferują taką funkcję. Zrzut ekranu pokazujący edycję tabeli w~tym programie pokazany jest na rysunku~\ref{rys:kile}.

%%%%%%%%%%%%%%%%%%%%%%%%%%%%%%%%%%%%%%%%%%%%%%%%
\paragraph{Wzory}
Jakość aplikacji do edycji tekstu najłatwiej poznać po typografii (estetyce) wzorów. W~tym obszarze prawdopodobnie nie ma lepszego narzędzia niż \LaTeX{}. Wzory zapisywane są za pomocą zbioru poleceń. Praktycznie tego samego zestawu znaczników do zapisu wzorów używa Wikipedia i~popularna biblioteka \href{https://www.mathjax.org/}{MathJax}.

Proste, bardziej popularne symbole matematyczne i~oznaczenia po pewnym czasie po prostu się pamięta. Jeśli chcemy uzyskać bardziej skomplikowane wzory lub użyć mniej popularnych symboli to warto sięgnąć po dokumentację dostępną na przykład na stronach \href{https://www.overleaf.com/learn/latex/Mathematical_expressions}{Overleaf -- Mathematical expressions}. Można też wspomóc się zewnętrzną aplikacją taką jak na przykład Kile. Alternatywnie można skorzystać z~edytorów równań dostępnych online, takich jak:
\begin{itemize}
	\item \href{https://latex.codecogs.com/eqneditor/editor.php}{CodeCogs}
	\item \href{https://www.latex4technics.com/}{Latex4technics}
\end{itemize}
Ich zaletą jest niemal natychmiastowe pokazywanie efektu bez konieczności czasochłonnej rekompilacji całego dokumentu.

Warto zwrócić uwagę na programy do detekcji symboli matematycznych:
\begin{itemize}
	\item \href{https://detexify.kirelabs.org/classify.html}{Detexify} -- narzędzie online,
	\item \href{https://github.com/zoeyfyi/TeX-Match}{TeX Match} -- wieloplatformowa, desktopowa wersja Detexify.
\end{itemize}

Przykłady wzorów matematycznych poniżej.

\begin{equation}
	R = \frac{U}{I}
\end{equation}

Wzory można też umieszczać w~tekście, otaczając je znakami dolara. Na przykład: $I=\frac{U}{R}$.

Kolejny wzór ma gwiazdkę przy \texttt{equation} i~dlatego nie będzie numerowany.

\begin{equation*}
	\vec{J} = nq \vec{v_d} =  \frac{nq^2 \tau}{m} \vec{E} = nq \mu \vec{E}    
\end{equation*}

Do wzoru można się odwoływać, jeśli ma on identyfikator i~numer. Wzór~\ref{eq:ohm} wiąże gęstość prądu $\pmb{J}$ z~natężeniem pola elektrycznego $\pmb{E}$ w~przewodniku.

\begin{equation}
	\pmb{J} = \sigma \pmb{E}
	\label{eq:ohm}
\end{equation}


Wzór Ampèra:

\begin{equation}
	\oint \vec{H} d\vec{l} = \int\limits_{S} \vec{J} \cdot d \vec{a} = I
\end{equation}

\begin{equation}
	\nabla \times \vec{H} = \vec{J}
\end{equation}


Twierdzenie Faradaya:

\begin{equation*}
	\Phi_B = \iint\limits_{\sum (t)} \pmb{B}(t) \cdot d\pmb{A}
\end{equation*}

\begin{equation}
	\mathcal{E} = -\frac{d\Phi_B}{dt}
\end{equation}

\begin{equation}
	\nabla \times \pmb{E} = - \frac{\partial \pmb{B}}{\partial t}
\end{equation}

Jeszcze kilka przypadkowo wybranych wzorów:

\begin{equation*}
	r = \sqrt{r_x^2 + r_y^2}
\end{equation*}

\begin{equation}
	\frac{dx}{dx} = 1
\end{equation}

\begin{equation}
	\frac{d}{dx} \ln x = \frac{1}{x}
\end{equation}

\begin{equation}
	y = \varphi  (y^\prime)x + \psi(y^\prime)
\end{equation}

\begin{equation}
	f(t) = \frac{1}{\pi} \int\limits_0^\infty \mathrm{d} \omega \,\int\limits_{-\infty}^{+\infty} f(\tau)\, \cos \omega (t-\tau)  \mathrm{d}\tau
\end{equation}

\[
\mathfrak{F}^{-1}\left(\mathfrak{F} \left[ f\left(t \right ) \right ] \right ) = f \left(t \right )
\qquad \textup{oraz} \qquad
\mathfrak{F}\left(\mathfrak{F}^{-1} \left[ F\left(j\omega \right ) \right ] \right ) = F \left(j\omega \right )
\]

\begin{equation*}
	x \oplus y = \neg (x \equiv y) = (x \vee y) \wedge (\neg x \vee \neg y) = (x \wedge \neg y) \vee (\neg x \wedge y)
\end{equation*}

\begin{equation*}
	\sum \frac{Q}{T} = 0
\end{equation*}

\begin{equation}
	\oint \frac{dQ}{T} = 0
\end{equation}

\begin{equation}
	\pmb{a} = \lim\limits_{\Delta t \to 0} \frac{\Delta \pmb{v}}{\Delta t} = \frac{d\pmb{v}}{dt}
\end{equation}

\begin{equation}
	K = m\,c^2 - m_0\,c^2 
	= m_0\,c^2 \left(
	\frac{1}{ \sqrt{1 - \left( \frac{v}{c} \right)^2 } } -1 
	\right)
\end{equation}

\begin{equation}
	\frac{\partial^2 y}{\partial x^2} = \frac{\mu}{F} \; \frac{\partial^2 y}{\partial t^2}
\end{equation}

\begin{equation}
	e^{\pm i\theta} = \cos \theta \pm i~\sin\theta
\end{equation}

\begin{equation}
	U(M) = \lim\limits_{R\to 0} \iiint\limits_{G-K}\frac{\rho(P)}{r}dG
\end{equation}

\begin{equation*}
	\left| \, \int\limits_{C_R} \frac{dz}{\left( 1+z^2 \right )^3} \, \right| \leqslant \sup\limits_{z \in C_R} \frac{1}{\left| 1+z^2 \right|^3} \; \pi R
\end{equation*}

W otoczeniu \texttt{align} można umieścić kilka wzorów jeden pod drugim oddzielając je znakiem nowego wiersza, czyli podwójnym ukośnikiem: \texttt{\textbackslash{}\textbackslash{}}. Wzory te poukładają się elegancko jeden pod drugim, jeśli w~każdej linii umieścimy jeden lub więcej znaków \&. Ważne, żeby w~każdej linii znak \texttt{\&} wystąpił tyle samo razy, gdyż wskazuje on miejsce wyrównania równań w~kolumnie. Przykłady poniżej.

\begin{align}
	g & = \frac{v^2}{R+h} \\
	v & = \sqrt{ \left( R+h \right) g }
\end{align}

\begin{align}
	y(t) = A_0 
	+& A_1 \sin \omega t + 
	A_2 \sin 2 \omega t + 
	A_3 \sin 3 \omega t + \ldots + \nonumber \\
	+& B_1 \cos \omega t + 
	B_2 \cos 2 \omega t + 
	B_1 \cos 3 \omega t + \ldots
\end{align}

\begin{align}
	\nu^\prime  &= \frac{\left(vt / \lambda \right) + \left( v_0t / \lambda \right)}{t} = \nonumber\\
	&= \frac{v + v_0}{\lambda} = \nonumber\\
	&= \frac{v + v_0}{v / \nu} = \nonumber\\
	&= \nu \frac{v+v_0}{v} = \nonumber\\
	&= \nu \left( 1 + \frac{v_0}{v} \right)
\end{align}

Nieco bardziej skomplikowane wzory z~eleganckim ułożeniem znaków równości w~pionie:

\begin{align}
	\oiint\nolimits_{\partial \Omega} \pmb{E} \cdot \mathrm{d}\pmb{S} &= \frac{1}{\varepsilon_0} \iiint\nolimits_\Omega \rho \, \mathrm{d}V &  \nabla \cdot \pmb{E} &= \frac {\rho} {\varepsilon_0} \label{eq:maxwell1} \\
	\oiint\nolimits_{\partial \Omega} \pmb{B} \cdot \mathrm{d}\pmb{S} &= 0 & \nabla \cdot \pmb{B} &= 0 \label{eq:maxwell2} \\
	\oint\nolimits_{\partial \Sigma} \pmb{E} \cdot \mathrm{d}\boldsymbol{\ell} &= -\frac{\mathrm{d}}{\mathrm{d}t}\iint\nolimits_{\Sigma}\pmb{B}\cdot\mathrm{d}\pmb{S} & \nabla \times \pmb{E} &= -\frac{\partial \pmb{B}}{\partial t} \label{eq:maxwell3} \\
	\oint\nolimits_{\partial \Sigma} \pmb{B} \cdot \mathrm{d}\boldsymbol{\ell} &= \mu_0 \iint\nolimits_{\Sigma} \pmb{J} \cdot \mathrm{d}\pmb{S} + \mu_0 \varepsilon_0 \frac{\mathrm{d}}{\mathrm{d}t} \iint\nolimits_{\Sigma} \pmb{E} \cdot \mathrm{d}\pmb{S} & \nabla \times \pmb{B} &= \mu_0\left(\pmb{J} + \varepsilon_0 \frac{\partial \pmb{E}} {\partial t} \right) \label{eq:maxwell4}
\end{align}

Otoczenie \texttt{gather} pozwala na łączenie na przykład macierzy ze ,,zwykłymi'' wzorami.

\begin{gather}
	\begin{bmatrix} \Phi_{11} & \Phi_{12} \\ \Phi_{21} & \Phi_{22} \end{bmatrix}
	=
	\frac{1}{\det(X)}
	\begin{bmatrix}
		X_{22} Y_{11} - X_{12} Y_{21} &
		X_{22} Y_{12} - X_{12} Y_{22} \\
		X_{11} Y_{21} - X_{21} Y_{11} &
		X_{11} Y_{22} - X_{21} Y_{12} 
	\end{bmatrix}
\end{gather}

\begin{gather}
	a \times b = -b \times a~= 
	\begin{vmatrix}
		\pmb{i}    &     \pmb{j}      &      \pmb{k}    \\
		a_x        &     a_y          &      a_z        \\
		b_x        &     b_y          &      b_z        
	\end{vmatrix}
	= \left(a_y b_z - b_y a_z \right)\pmb{i} 
	+ \left(a_z b_x - b_z a_x \right)\pmb{j} 
	+ \left(a_x b_y - b_x a_y \right)\pmb{k} 
\end{gather}

% Na Wydziale Elektrycznym wzorów chemicznych raczej nie potrzebujemy zbyt często a~te kilka przypadków da się zapisać jako zwykłe równania.
% Gdyby jednak ktoś bardzo chciał to można odkomentować linijkę z~pliku cls:
%\RequirePackage[version=4]{mhchem}
% a~wtedy zadziałają poniższe makra:
%\ce{3H2O} \\
%\ce{1/2H2O} \\
%\ce{AgCl2-} \\
%\ce{H2_{(aq)}}

%%%%%%%%%%%%%%%%%%%%%%%%%%%%%%%%%%%%%%%%%%%%%%%
\paragraph{Kody źródłowe}
Do umieszczania kodów źródłowych w~szablonie wykorzystany jest pakiet \texttt{listings}. Została przygotowana propozycja wyglądu listingów oparta o~prezentowane wcześniej kolory. Ustawienia można zmienić w~pliku \texttt{cls}. Poniżej znajduje się kilka przykładów wykorzystania wspomnianego pakietu. Więcej informacji można znaleźć na przykład w~dokumentacji Overleaf\footnote{\href{https://www.overleaf.com/learn/latex/code_listing}{<Overleaf -- Code listing>}}.

Wklejanie kodu w~postaci zrzutu ekranu (rysunku) jest niedopuszczalne gdyż \textit{tekst} kodu źródłowego jest \textit{tekstem}, nawet jeśli listing jest traktowany jako element graficzny.

Kod źródłowy można umieścić bezpośrednio w~pliku tex, tak jak poniższy przykład z~listingu~\ref{lst:hellopy}. Jednak z~wielu powodów nie jest dobry pomysł.

\begin{lstlisting}[language=Python,
	caption={Prosty skrypt w~języku Python},
	label={lst:hellopy}]
	#!/usr/bin/env python
	# -*- coding: utf-8 -*-
	"""Simple world of hello.
	"""
	
	import sys
	
	def main():
		"""The one and only function"""
		fib = lambda n: reduce(lambda x, n: [x[1], x[0]+x[1]], range(n), [0, 1])[0]
		try:
			print(fib(int(sys.argv[1])))
		except:
			print("Hello World!")
	
	if __name__ == "__main__":
		main()
\end{lstlisting}

Kolorowanie składni i~numerowanie linii dzieją się automatycznie przy czym tekst kodu nadal jest tekstem.

Spójna kolorystyka kodów w~różnych językach może być pożądana albo wręcz przeciwnie. Dlatego można rozważyć jakąś, niewielką, lokalną zmianę stylu tak, jak w~poniższym przykładzie. Przykład kodu źródłowego w~innym języku (tutaj: C) jest pokazany w~listingu~\ref{lst:helloC}. 

\begin{lstlisting}[language=C,
	backgroundcolor=\color{bezowy!25!white},
	caption={Prosty kod w~języku C},
	label={lst:helloC}]
	#include <stdlib.h>
	#include <stdio.h>
	
	int main(int argc, char **argv) {
		printf("Hello World!\n");
		return EXIT_SUCCESS;
	}
\end{lstlisting}

Zauważ bardzo niefortunne przejście na nową stronę, dodatkowo zakłócone przez przypis dolny, czego niestety \LaTeX{} nie jest w~stanie samodzielnie poprawić. Należałoby powyżej dopisać więcej tekstu by ,,przepchnąć'' kod na kolejną stronę, albo tę ilość tekstu zmniejszyć albo podzielić listing na mniejsze fragmenty i~omówić osobno, albo przesunąć w~treści -- działania do podjęcia są prawie takie same, jak w~przypadku rysunków.

Kod źródłowy lepiej trzymać w~zewnętrznym pliku i~załączać go do dokumentu dynamicznie, tak jak ma to miejsce w~przypadku listingu~\ref{lst:Keccak}. Można wskazać zakres linii, które powinny pojawić się w~dokumencie. Jednak wtedy zapewne warto pamiętać o~ustawieniu numeru pierwszej linii aby numeracja w~dokumencie zgadzała się z~faktycznym numerem linii w~pliku. Listing~\ref{lst:Keccak} ma automatycznie ustawiany podpis na podstawie nazwy pliku. W~pewnych sytuacjach może być to wygodniejsze niż ręczne wprowadzanie nazwy.

\lstinputlisting[language=C,
firstnumber=31,
firstline=31,
lastline=50,
caption=\lstname,
label={lst:Keccak}]{Keccak-inplace32BI.c}

W otoczeniu listing nie można stosować polskich znaków diakrytycznych.

%%%%%%%%%%%%%%%%%%%%%%%%%%%%%%%%%%%%%%%%%%%%%%%
\paragraph{Ozdobniki graficzne w~opisie oprogramowania}
Szablon wczytuje pakiet pozwalający na wyróżnienie w~tekście informacji o~skrótach klawiszowych, poruszaniu się po menu programu i~ścieżki plików. Oczywiście nie ma konieczności korzystania z~tego dodatku, jeśli ktoś nie uważa go za potrzebny.

W tekście można wyróżnić skróty klawiszowe, takie jak na przykład: \keys{\Alt, F4}, \keys{\ctrl, \Alt, \del}.
\begin{itemize}
	\item \keys{A, a, B, b, C, c, 1, 2, 3, PgUp}
	\item \keys{\Space} \keys{\SPACE}
	\item \keys{\backspace} \keys{\del} \keys{\backdel}
	\item \keys{\return} \keys{\enter}
	\item \keys{\shift} \keys{\capslock}
	\item \keys{\ctrl} \keys{\Alt} \keys{\AltGr}
	\item \keys{\tab}
	\item \keys{\esc} \keys{\oldesc}
	\item \keys{\winmenu}
	\item \keys{\arrowkey{^}} \keys{\arrowkeyup}
	\item \keys{\arrowkey{v}} \keys{\arrowkeydown}
	\item \keys{\arrowkey{>}} \keys{\arrowkeyright}
	\item \keys{\arrowkey{<}} \keys{\arrowkeyleft}
\end{itemize}

Jeśli opisujemy poruszanie się po menu to również można użyć ozdobnika graficznego: \menu{View > Zoom > Zoom in}.

Podobne rozwiązanie jest dla ścieżek katalogów: \directory{/etc/apache2/mods-available}.

%%%%%%%%%%%%%%%%%%%%%%%%%%%%%%%%%%%%%%%%%%%%%%%
\paragraph{Akronimy i~symbole}
Zgodnie z~obowiązującym Zarządzeniem pod koniec pracy znajduje się lista symboli i~akronimów pod nazwą ,,Wykaz symboli i~skrótów''. Na liście wystąpią tylko te akronimy, które faktycznie zostaną użyte w~tekście, za pomocą jednej z~pokazanych niżej metod. Jeśli żaden akronim lub symbol w~tekście nie wystąpi, w~sensie użycia odwołania za pomocą odpowiedniego polecenia, to generowanie strony z~listą akronimów zostanie pominięte.

Szablon używa pakietu \texttt{glossaries} do zarządzania akronimami i~symbolami. Listę tych elementów należy przygotować w~pliku \texttt{glossary.tex}, zgodnie z~pokazanym tam szablonem. Aby w~pracy pojawiła się lista akronimów należy użyć polecenia:

\begin{lstlisting}[language=bash,
	numbers=none,
	caption=Wygenerowanie listy skrótów i~symboli,
	label={lst:gloss}]
	makeglossaries [nazwa pliku podstawowego bez rozszerzenia tex]
\end{lstlisting}

Pierwsze użycie akronimu w~tekście jest rozpoznawane automatycznie i~pojawia się on w~pełnej oraz skróconej formie: \gls{CPEE}. Kolejnymi razy prezentuje się już tylko wersja skrócona, chociaż oba wywołania w~\LaTeX{u} wyglądają tak samo: \gls{CPEE}. Dzięki temu nie trzeba się zastanawiać, czy dany akronim został wcześniej pokazany w~pełnej formie, czy też jeszcze nie.

Można też samodzielnie wybrać jaką formę ma przybrać akronim w~tekście, co nie wpływa na rozpoznawanie pierwszego użycia skrótu:
\begin{itemize}
	\item krótkiej: \acrshort{IEEE}, \acrshort{PW}, \acrshort{IETiSIP},
	\item długiej: \acrlong{IEEE}, \acrlong{PW}, \acrlong{IETiSIP},
	\item pełnej: \acrfull{IEEE}, \acrfull{PW}, \acrfull{IETiSIP}.
\end{itemize}

Chociaż na powyższej liście występuje \gls{PW} to tutaj akronim jest traktowany jak pojawiający się po raz pierwszy. Dopiero kolejne użycie daje inny efekt: \gls{PW}.

W Overleaf wykorzystanie wykazu symboli matematycznych okazuje się sprawiać problemy -- indeks nie odświeża się automatycznie wtedy, gdy powinien. Żeby zmusić Overleaf do odświeżenia listy symboli i~akronimów, należy z~rozwijanego menu przy zielonym przycisku ,,Recompile'' wybrać opcję ,,Recompile from scratch''. Trwa to nieco dłużej ale spełnia swoją rolę.

Przykład użycia symbolu natężenia prądu elektrycznego~\gls{symb:I}. Jedną z~ciekawszych liczb jest~\gls{symb:Pi}. Tak użyte symbole przy prawidłowej kompilacji pliku \texttt{PDF} znajdą się na indeksowanej liście a~więc będzie można znaleźć miejsca ich użycia w~pracy. Oczywiście to indeksowanie zadziała tylko wtedy, gdy odwołamy się do symbolu w~przedstawiony sposób a~nie na przykład za pomocą otoczenia matematycznego: $\pi$.

%%%%%%%%%%%%%%%%%%%%%%%%%%%%%%%%%%%%%%%%%%%%%%%
\paragraph{Cytowanie a~cytaty}
Cytaty to nie to samo, co cytowania i~to nie to samo co odnośniki bibliograficzne. Cytaty nie są popularne w~naukach technicznych. Jednak przywołując czyjeś słowa trzeba cytat stosownie oznaczyć za pomocą cudzysłowów i~podając autora.

Cytując \textbf{nie należy} korzystać z~dostępnego na klawiaturze znaku \keys{"{}}. Taki znak nie jest zgodny z~polską typografią. Cytowania tekstu standardowo robi się ,,recznie'', oznaczając cudzysłowy za pomocą podwójnych przecinków i~apostrofów. Wiele programów (np.: TeXstudio) potrafi automatycznie zmieniać cudzysłów wybrany z~klawiatury w~polskie znaki cytowania -- jest to opcja do włączenia w~ustawieniach i~bardzo ułatwia pisanie.

Dodatki usprawniające obsługę cytowań są dostępne dla \LaTeX{a}, ale zapewne na uczelni technicznej nie ma potrzeby ich stosowania.

%%%%%%%%%%%%%%%%%%%%%%%%%%%%%%%%%%%%%%%%%%%%%%%
\paragraph{Kompilowanie lokalnie}
Darmowe konto w~Overleaf zapewnia czas każdej kompilacji nie dłuższy niż 1 minuta. W~przypadku skomplikowanej pracy czas ten może zostać przekroczony i~wówczas Overleaf nie wygeneruje pliku wynikowego. Dlatego warto rozważyć pobranie szablonu na własny komputer z~zainstalowanym \LaTeX{em}. Kompilacja może wówczas wyglądać tak, jak na listingu~\ref{lst:kompilacja}.

\begin{lstlisting}[language=bash,
	caption={Kompilacja pracy dyplomowej lokalnie},
	label={lst:kompilacja}]
	pdflatex EE-dyplom && biber EE-dyplom && makeglossaries EE-dyplom && pdflatex EE-dyplom && pdflatex EE-dyplom
\end{lstlisting}

Polecenie \textbf{pdflatex} można zastąpić przez \textbf{xelatex} i~dla obu kompilatorów rezultat wynikowy powinien być niemal identyczny. Jednak nie należy mieszać obu rozwiązań gdyż ich pliki pośrednie zakłócają się nawzajem w~procesie kompilacji.

Do edycji i~kompilacji można skorzystać z~dedykowanych aplikacji takich jak:
\begin{itemize}
	\item \href{https://miktex.org/}{MiKTeX} -- dla Windows, dostępny w~MS Store, zawierający niezbędne narzędzia \LaTeX{a} i~prosty edytor \href{https://www.tug.org/texworks/}{TeXworks}.
	\item \href{https://www.texstudio.org/}{TeXstudio} -- środowisko bardziej rozbudowane niż TeXworks ale łatwe w~użyciu, potrzebuje kompilatora \LaTeX{a} więc pod Windows można je zainstalować po przetestowaniu MiKTeXa.
	\item \href{https://kile.sourceforge.io/}{Kile} -- podobne do TeXstudio.
\end{itemize}

Są to aplikacje wieloplatformowe dostępne dla Linux, Windows, macOS. Użytkownicy systemu Linux najpewniej pobiorą odpowiednie oprogramowanie, zaczynając od ,,\TeX Live'', z~repozytoriów używanej dystrybucji.

\paragraph{Więcej informacji na temat \LaTeX{a}}
\begin{itemize}
	\item \href{https://www.overleaf.com/learn}{<https://www.overleaf.com/learn>} -- przystępny tutorial na stronie Overleaf,
	\item \href{https://www.latex-project.org/}{<https://www.latex-project.org/>} -- strona domowa projektu,
	\item \href{https://www.tug.org/begin.html}{<https://www.tug.org/begin.html>} -- dobry zbiór odnośników do innych materiałów.
\end{itemize}


\end{easyappendix}

%%%%%%%%%%%%%%%%%%%%%%%%%%%%%%%%%%%%%%%%%%%%%%%%%%%%%%%%%%%%%%%%%%%%%%%%%
\begin{easyappendix}{Dowód próżni doskonałej}
\end{easyappendix}

%%%%%%%%%%%%%%%%%%%%%%%%%%%%%%%%%%%%%%%%%%%%%%%%%%%%%%%%%%%%%%%%%%%%%%%%%
\begin{easyappendix}{Dowód nieskończoności urojonej}
\end{easyappendix}
